\documentclass[10pt,a4paper]{amsart}
\usepackage[margin=19mm]{geometry}
\usepackage{amsmath,amssymb,mathtools}
\usepackage[scr=rsfs,cal=euler]{mathalfa}
\usepackage{xfrac}
\usepackage[unicode,hidelinks,bookmarks=false]{hyperref}
\newcommand{\ie}{i.e.\ }
\newcommand{\cf}{cf.\ }
\newcommand{\resp}{respectively\ }
\newcommand{\Subsection}{\subsection}
\providecommand{\nobreakdash}{\nobreak\hbox{-}}
\setlength{\emergencystretch}{2em}
\multlinegap0pt
\usepackage{bm}
\usepackage{bbm}
\usepackage{dsfont}
\usepackage{xspace}
\usepackage{cancel}
\usepackage{mathtools}
\usepackage{amsthm}
\makeatletter
\@ifundefined{thm}{\newtheorem{thm}{Thm}}{}
\@ifundefined{lemma}{\newtheorem{lemma}[thm]{Lemma}}{}
\makeatother
\newcommand{\cS}{\mathcal S}
\newcommand{\ds}{{\mathrm ds}}
\newcommand{\dw}{{\mathrm dw}}
\newcommand{\io}{{\iota}}
\newcommand{\tS}{{\widetilde S}}
\title[Perfect t-Embeddings of Uniformly Weighted Generalized Tower]{Perfect t-Embeddings of Uniformly Weighted Generalized Tower Graphs}
\author{David Keating, Hieu Trung Vu}
\date{Source: arXiv:2509.18791v1 (CC BY 4.0)}
\begin{document}
\maketitle

\noindent\emph{Source and licence.} Perfect t-Embeddings of Uniformly Weighted Generalized Tower Graphs. Version arXiv:2509.18791v1 (CC BY 4.0), distributed under the Creative Commons Attribution 4.0 International licence, as recorded in the arXiv licence metadata for this exact version and shown on the abstract page https://arxiv.org/abs/2509.18791v1. Full rights notice: licenses/arxiv-2509.18791-CC-BY-4.0.txt.\par\smallskip

\noindent\textit{Editorial scope.} Counted unit: the Lemma whose source label is \texttt{lemma: steepest\_descent\_lemma} in the article (source0.tex lines 1460--1471, source label \texttt{lemma: steepest\_descent\_lemma}), delivered together with the article's own complete original proof of it (source0.tex lines 1474--1532, one \texttt{proof} environment of 59 lines). No mathematics is written, repaired or re-derived: statement and proof are the article's own text line for line. Required context retained uncounted: none. Every definition, display and label of those lines is retained unchanged. Omitted from this package and not redistributed: 1--1459, 1472--1473, 1533--1894. Nothing inside a retained excerpt is omitted or altered: every retained body is delivered exactly as the article's lines are written, so no omission receipt of any kind accompanies this package. Citations: the retained excerpts carry no citation command at all, so no bibliography entry of the article is transcribed and no reference is redistributed. Reference adaptation: none was needed and none was made; every reference inside the delivered unit resolves inside it, so no unresolved reference or citation is produced. Printed slips kept literal: no printed word, symbol, hypothesis or number of the counted statement or of its proof was altered. Layout and class: the article's own class is \texttt{amsart} with packages including graphicx, epstopdf, amsmath, xcolor, float, mathtools, amsthm, cite, stix, amsfonts, setspace, adjustbox, subcaption, tikzit; the wrapper is the lane's pinned \texttt{amsart} editorial wrapper at 10pt on A4, which restates the theorem-like environments the retained text needs and adds the article's own macro definitions that the retained text needs (\texttt{\textbackslash cS}, \texttt{\textbackslash ds}, \texttt{\textbackslash dw}, \texttt{\textbackslash io}, \texttt{\textbackslash tS}), with extra wrapper packages (bm, bbm, dsfont, xspace, cancel, mathtools). The delivered numbering is the unit's own. Assets: no retained span contains a figure environment, a graphics-inclusion command, a PDF-page inclusion, a table or a TikZ picture; no image file is copied into this package, the package declares no asset dependency and no image file is redistributed with it. Licence: version arXiv:2509.18791v1 of this article is distributed under the Creative Commons Attribution 4.0 International licence, as recorded in the arXiv licence metadata for this exact version and shown on the abstract page https://arxiv.org/abs/2509.18791v1. A full rights notice is delivered at licenses/arxiv-2509.18791-CC-BY-4.0.txt. Characters: every retained excerpt is pure ASCII, so the delivered engine, XeLaTeX with its default Latin Modern text font, reports no missing character.
\begin{lemma}
    \label{lemma: steepest_descent_lemma}
    \begin{equation}
      \label{eq: asymptotic of w integral 
}
\begin{aligned}
&\quad I_w(n):=\int_{\mathcal{E}_2}\mathrm e^{(n+1)\tS(w)}\frac{w-\mathrm i}{w}\left (\frac{w}{(1+w)(w-1)}\right )^{\frac{\epsilon(j,k)+1}{2}} G_{\io,\io'}(w)\operatorname{d}w\\
        &\sim G(\xi)\mathrm e^{(n+1) \tS(\xi)}\sqrt{\frac{\pi}{n+1}}\sqrt{\frac{2}{|\tS''(\xi)|}}e^{-i\theta_\xi}+G(\bar \xi)e^{(n+1) \tS(\bar \xi)}\sqrt{\frac{\pi}{n+1}}\sqrt{\frac{2}{|\tS''(\bar \xi)|}}e^{-i\theta_{\bar \xi}}
\end{aligned}
    \end{equation}
where $$G(w):=\frac{w-\mathrm i}{w}\left (\frac{w}{(1+w)(w-1)}\right )^{\frac{\epsilon(j,k)+1}{2}}  G_{\io,\io'}(w).$$
\end{lemma}

 \begin{proof}
     We perform the steepest descent asymptotic analysis on the integral $I_w(m)$. First, recall that the action $\tS(z)$ has two critical points/saddle points that are simple and also conjugate pairs.
     And since $\tS(z)$ is analytic in the upper half plane $\mathbb H$, the paths of steepest ascent and descent are the level lines
of $\operatorname{Im}(\tS(z))=\operatorname{Im}(\tS(\xi))$. These paths can only cross the real line at 0. We now describe
these paths. For the saddle point $\xi$ on the upper-half plane, the steepest ascent contour start from $0$, passes through $\xi$ and goes to $\infty$. The complete steepest ascent includes the reflection of the mentioned path in the lower half plane. We will denote the steepest ascent contour in $w$ as $\gamma_w$. Note that the contour $\mathcal E_2$ only has a singularity at $w=1$. Thus, we can deform $\mathcal E_2$ to the steepest descent contour without picking up the residue at $w=1$. The Taylor expansion near the saddle point $\xi$ of the action $\tS(w)$ gives,
     \[
     \tS(w) = \tS(\xi)+\frac{1}{2}\tS''(\xi)(w-\xi)^2 +O(w^3)
     \]
     since the first derivative vanishes at $\xi$. Writing $ \tS(w)= u(w)+\mathrm i  v(w)$,  then along the steepest descending path, the integral $I_\xi(n)$ is approximated, 
     \begin{equation}
         \label{eq: I(n) approx}
    \begin{aligned}
    I_\xi(n)&=\mathrm e^{(n+1)\tS(\xi)}\int_{\mathcal E_2}G(w)\mathrm e^{(n+1)(\tS(w)-\tS(\xi))}\operatorname{d}w\\
    &\sim G(\xi)\mathrm e^{(n+1)\tS(\xi)}\int_{\mathcal E_2}\mathrm e^{\mathrm i (n+1)(v(w)-v(\xi))}\mathrm e^{(n+1)(u(w)-u(\xi))}\operatorname{d}w \\
    &\sim G(\xi)\mathrm e^{(n+1) \tS(\xi)}\int_{\gamma_w}\mathrm e^{(n+1)(u(w)-u(\xi))}\mathrm dw
    \end{aligned}    
     \end{equation}
     where $\gamma_w$ is the steepest descent contour. Note that along this contour, the imaginary part of $\widetilde\cS$ stays constant so $v(w)=v(\xi)$. Let $s^2=u(\xi)-u(w)$, then we can rewrite \eqref{eq: I(n) approx} as
    \begin{equation}
        \label{eq: I(n) approx 2}
        \begin{aligned}
            I_\xi(n) \sim G(\xi)\mathrm e^{(n+1)\tS(\xi)}\int_{-\infty}^\infty \mathrm e^{-(n+1)s^2}\left(\frac{\mathrm dw}{\mathrm ds}\right)\mathrm ds.
        \end{aligned}
    \end{equation}
    Note that we have the Gaussian integral $\displaystyle \int_{-\infty}^\infty \mathrm e^{-(n+1)s^2} \ds=\sqrt{\frac{\pi}{n+1}}$, so the leading term of the asymptotic expression near the saddle point where $s=0$ gives 
    \[
    I_\xi(n) \sim G(\xi)\mathrm e^{(n+1)\tS(\xi)}\left(\frac{\mathrm dw}{\mathrm ds}\right)_{s=0} \sqrt{\frac{\pi}{n+1}}.
    \]
    On the other hand, since $v(w)$ is constant along the steepest descent path, $-s^2=u(w)-u(\xi)=\tS(w)-\tS(\xi)$. By taking derivative with respect of $w$ twice, we have 
    \[-2s\left(\frac{\ds}{\dw}\right)=\tS'(w),\]
    \[
    -2\left(\frac{\ds}{\dw}\right)\left(\frac{\ds}{\dw}\right)-2s\frac{\mathrm d^2s}{\dw^2}=\tS''(w)
    \]
    and so $\displaystyle \left(\frac{\dw}{\ds}\right)_{s=0}=\left(\frac{-2}{\tS''(\xi)}\right)^{1/2}$. 
   This square root has some ambiguity in sign. To avoid this issue, denote $\displaystyle \theta_{SDP}:=\operatorname{arg}\left(\frac{\mathrm d w}{\mathrm d s}\right)_{s=0}$. Then, 
  $ 
   \displaystyle \left(\frac{\dw}{\ds}\right)_{s=0}=\left(\frac{2}{|\tS''(\xi)|}\right)^{1/2}e^{i\theta_{SDP}}$. 
   On the other hand, recall that we define $\theta_\xi:=\frac{1}{2}\operatorname{arg}(\tS''(\xi))$. Let $\tS''(\xi)=R\mathrm e^{\mathrm 2i\theta_\xi}$ and $w-\xi=r\mathrm e^{\mathrm i \theta_{SDP}}$, we have,
   \[
     \tS(w)-\tS(\xi) = \frac{1}{2}\tS''(\xi)(w-\xi)^2 +O(w^3)\sim\frac{1}{2}(Rr^2)\mathrm e^{2i(\theta_\xi+\theta_{SDP})}.
     \]
    But note that since we are on the steepest descent path $\gamma_w$, $2(\theta_\xi+\theta_{SDP})=\pm \pi$ and thus $\theta_{SDP}=-\theta_\xi\pm\frac{\pi}{2}$ (where the sign is dependent on the direction of the contour) and we have the asymptotic of $I_\xi(n)$ given by, 
    \begin{equation}
        \label{eq: I(n) approx 3}
        I_\xi(n)\sim G(\xi)\mathrm e^{(n+1) S(\xi)}\sqrt{\frac{\pi}{n+1}}\sqrt{\frac{2}{|S''(\xi)|}}\mathrm e^{-i\theta_\xi}.
    \end{equation}
 Recall that the contour $\mathcal E_2$ has also passes through the other critical point $\bar \xi$ on the lower half-plane. The analysis is similar and we have
 \begin{equation}
        \label{eq: I(m) approx 4 conjugate}
        I_{\bar \xi}(n)\sim G(\bar \xi)\mathrm e^{(n+1)S(\bar \xi)}\sqrt{\frac{\pi}{n+1}}\sqrt{\frac{2}{|S''(\bar \xi)|}}\mathrm e^{-i\theta_{\bar \xi}}.
    \end{equation}
Altogether, we have that the leading order asymptotic of the integral $I_w(n)$ is given by
$$
\begin{aligned}
I_w(n)&\;=I_\xi(n)+I_{\bar \xi}(n)\\
&\;\sim G(\xi)\mathrm e^{(n+1) \tS(\xi)}\sqrt{\frac{\pi}{n+1}}\sqrt{\frac{2}{|\tS''(\xi)|}}\mathrm e^{-i\theta_\xi}+G(\bar \xi)\mathrm e^{(n+1) \tS(\bar \xi)}\sqrt{\frac{\pi}{n+1}}\sqrt{\frac{2}{|\tS''(\bar \xi)|}}\mathrm e^{-i\theta_{\bar \xi}}.
\end{aligned}
$$
 \end{proof}

\end{document}
